%% Chapter 1: Electrostatics  (textbook Topics 15--16)
%% Compilable on its own:  make latex CH=ch01-electrostatics
\documentclass[../book.tex]{subfiles}
% Resolve (and link) references into other chapters when the book has been
% built in this directory; latexmk keeps its .aux in build/, see .latexmkrc.
% Ignored entirely by the whole-book build, and a silent no-op if book.aux is
% absent -- so a stale path here costs cross-references with no error.
\IfFileExists{../build/book.aux}{\externaldocument{../build/book}}{}

\begin{document}

\chapter{Electrostatics}
\label{ch:electrostatics}

\section{Coulomb's law}

\begin{figure}[htbp]
  \centering
  \tikzfig{ch01-coulomb-law.tex}
  \caption{Coulomb's law. The source charge $Q$ sits at $\mathbf{r}'$ and the
  charge $q$ sits at $\mathbf{r}$, which feels the force $\mathbf{F}$ due to $Q$.}
  \label{fig:coulomb-law}
\end{figure}


\begin{keyeq}
  \label{eq:coulomb-law}
  \mathbf{F}
  =
  k_e 
  \frac{Q \, q}{\abs{\mathbf{R}}^2}
  \hat{\mathbf{R}}
  \;,
\end{keyeq}
where
\begin{align}
    k_e
    &=
    \frac{1}{4 \pi \epsilon_0}
    \approx 9 \times 10^9 \;\mathrm{N \cdot m^2 / C^2}
    \\
    \mathbf{R}
    &=
    \mathbf{r} - \mathbf{r}'
    \\
    \hat{\mathbf{R}}
    &=
    \frac{
    \mathbf{R}
    }{
    \abs{\mathbf{R}}
    }
    \;.
\end{align}



\section{The electric field}
Instead of having a Coulomb force between two charges $Q$ and $q$, we can have a different interpretation: The charge $Q$ does not interact with charge $q$ directly. Instead, the charge $Q$ creates an \emph{electric field} $\mathbf{E}(\mathbf{r})$ around it, and the charge $q$ at position $\mathbf{r}$ feels a force from the electric field $\mathbf{E}(\mathbf{r})$ locally at $\mathbf{r}$. 

The electric field $\mathbf{E}$ is a \emph{vector field}. We may think of a vector field as: ``at every point in the space $\mathbf{r}$, there is a vector $\mathbf{V}(\mathbf{r})$ (a little arrow) attached to that point $\mathbf{r}$.'' Because of that, a vector field $\mathbf{V}(\mathbf{r})$ is a function of location, i.e., space coordinate, $\mathbf{r}=(x,y,z)$.

\begin{figure}[htbp]
  \centering
  \tikzfig{ch01-coulomb-law.tex}
  \caption{Same as \autoref{fig:coulomb-law}}
  \label{fig:coulomb-law-E-field}
\end{figure}

To make our two-charge example more explicit, we can refer to \autoref{fig:coulomb-law-E-field}. Here, charge $Q$ creates an electric field:
\begin{align}
    \label{eq:coulomb-law-E-field}
    \mathbf{E}(\mathbf{r})
    &=
    k_e 
    \frac{Q}{\abs{\mathbf{R}}^2}
    \hat{\mathbf{R}}
    \\
    &=
    k_e 
    \frac{Q}{\abs{\mathbf{r} - \mathbf{r}'}^2}
    \frac{\mathbf{r} - \mathbf{r}'}{\abs{\mathbf{r} - \mathbf{r}'}}
    \;,
\end{align}
and change $q$ feels the force
\begin{align}
    \mathbf{F}(\mathbf{r})
    &=
    q \mathbf{E}(\mathbf{r})
    \\
    &=
    k_e 
    \frac{Q \, q}{\abs{\mathbf{R}}^2}
    \hat{\mathbf{R}}
    \;.
\end{align}
Quantitatively, this is the same as the Coulomb force in the static case, and it seems the $\mathbf{E}$ field is just a mathematical tool to facilitate calculation. However, this is in fact a deep breakthrough in theory, which later on leads to the prediction of the existence of electromagnetic waves.


\section{Gauss's law}


\section{Electric potential energy}


\section{Capacitance and dielectrics}


\chapterbib

\end{document}
